\begin{lemma}
Let $1<p<\infty$ and let $\gamma_n=\gamma_n(p)$ be as in \noderef{GammaExp}.
Then
\begin{enumerate}
\item[(a)] $\gamma_n>0$ for every integer $n\ge1$, and
\item[(b)] $\dfrac{1}{\gamma_1}\le\dfrac{n}{\gamma_n}$ for every integer
$n\ge1$.
\end{enumerate}
That is, the smallest of the numbers $n/\gamma_n$, $n=1,2,\dots$, occurs for
$n=1$, as asserted in Section~6 of the paper.
\end{lemma}
\begin{proof}
Fix $p>1$, so $p-1>0$.  By \noderef{GammaExp} and \noderef{LambdaCoef}, for
every integer $n\ge1$,
\[
  2\gamma_n=2\lambda_{n+1}(n,p)=-np+\sqrt{4(n+1)^2(p-1)+n^2(p-2)^2},
\]
and $np>0$.  Write $R(n)=4(n+1)^2(p-1)+n^2(p-2)^2$, a sum of two nonnegative
terms, so $R(n)\ge0$.

\emph{Step 1 (proof of (a)).}  Fix $n\ge1$.  Then
\[
  R(n)-n^2p^2=4(n+1)^2(p-1)+n^2\bigl[(p-2)^2-p^2\bigr]
  =4(n+1)^2(p-1)-4n^2(p-1)=4(p-1)(2n+1),
\]
where we used the identity $(p-2)^2-p^2=-4(p-1)$ and
$(n+1)^2-n^2=2n+1$.  Since $p-1>0$ and $n\ge1$, this is positive.  Hence
$R(n)>n^2p^2\ge0$, and therefore
\[
  \sqrt{R(n)}>\sqrt{n^2p^2}=np
\]
(the square root being strictly increasing on $[0,\infty)$, and $np\ge0$).
Consequently $2\gamma_n=-np+\sqrt{R(n)}>0$, i.e.\ $\gamma_n>0$.

\emph{Step 2 (a scale-free formula for $\gamma_n/n$).}  Fix $n\ge1$.  Since
$n>0$,
\[
  \frac{\sqrt{R(n)}}{n}=\sqrt{\frac{R(n)}{n^2}}
  =\sqrt{4\Bigl(\frac{n+1}{n}\Bigr)^2(p-1)+(p-2)^2}
  =\sqrt{4\Bigl(1+\frac1n\Bigr)^2(p-1)+(p-2)^2},
\]
and therefore
\[
  \frac{2\gamma_n}{n}=-p+\sqrt{4\Bigl(1+\frac1n\Bigr)^2(p-1)+(p-2)^2}.
\]

\emph{Step 3 ($\gamma_n/n\le\gamma_1$).}  Fix $n\ge1$.  Then $1/n\le1$, so
$0<1+\frac1n\le2$ and hence $\bigl(1+\frac1n\bigr)^2\le4$.  Since $p-1>0$,
\[
  4\Bigl(1+\frac1n\Bigr)^2(p-1)+(p-2)^2\ \le\ 16(p-1)+(p-2)^2 ,
\]
and the square root is nondecreasing on $[0,\infty)$, so
\[
  \frac{2\gamma_n}{n}=-p+\sqrt{4\Bigl(1+\tfrac1n\Bigr)^2(p-1)+(p-2)^2}
  \ \le\ -p+\sqrt{16(p-1)+(p-2)^2}=\frac{2\gamma_1}{1}=2\gamma_1 ,
\]
the right-hand equality being Step 2 at $n=1$, where
$4\bigl(1+\frac11\bigr)^2=16$.  Dividing by $2$ gives
$\gamma_n/n\le\gamma_1$.

\emph{Step 4 (proof of (b)).}  Fix $n\ge1$.  By Step 1 both $\gamma_n$ and
$\gamma_1$ are positive, and $n>0$, so $\gamma_n/n>0$.  Step 3 gives
$0<\gamma_n/n\le\gamma_1$, and taking reciprocals of positive numbers reverses
the inequality:
\[
  \frac{1}{\gamma_1}\le\frac{1}{\gamma_n/n}=\frac{n}{\gamma_n}.
\]
\end{proof}
